
\section{Verification of the algebra}
\label{sec:alg}

The relations among the one-loop components are the arity-two components of $L\circ L=0$; they hold on the cohomology of $\Lambda^{(0)}$ as soon as the quantum master equation holds through arity two. Without restoring counterterms they fail by the unrestored obstruction, which makes each relation a computable test of the scheme.

\begin{theorem}[Relations and their defect]
\label{thm:rel}
Let $S(\Gamma_{\rm ren})=\hbar\,\mathcal B+O(\hbar^2)$ with $\mathcal B:=\mathrm{Ob}-sF_1$. Then $S_\Gamma^2=\hbar\,\{\mathcal B,-\}+O(\hbar^2)$ on renormalised insertions, and on local operators
\[
L\circ L=\hbar\,\mathrm{ad}_{\mathcal B}+O(\hbar^2).
\]
With $\delta:=$ the order-$\hbar$ part of $-\Lambda^{(1)}$ on cohomology and $\rho_0$ the classical action [HK, def:c12-class], the arity-two component at order $\hbar$ on cocycles $O$ reads
\begin{equation}
\begin{aligned}
[\delta(x),\rho_0(y)]+[\rho_0(x),\delta(y)]-\delta([x,y])={}&\pm\{\mathcal B^{(2)}(x,y),-\}\pm\big[\{\mathcal B^{(1)}(x),-\},\rho_0(y)\big]\\
&\pm\big[\rho_0(x),\{\mathcal B^{(1)}(y),-\}\big]+\Lambda^{(0)}\text{-exact},
\end{aligned}
\label{eq:reldefect}
\end{equation}
with the signs of the cochain bracket [HK, def:c2-dce]. If $sF_1=\mathrm{Ob}$ through arity two along $(x,y)$, the right side is $\Lambda^{(0)}$-exact and the relations (R2)--(R8) of [HK, cor:c12-rel] hold on cohomology.
\end{theorem}

\begin{proof}
For any functional $\Gamma$, $S_\Gamma^2=(S(\Gamma),-)$ by the Jacobi identity and the derivation property of $d_{CE}$. Applied to $[M]_R$ and \eqref{eq:repD}, $[L(L(M))]_R=S_\Gamma^2[M]_R=\hbar(\mathcal B,[M]_R)+O(\hbar^2)$, and by Theorem \ref{thm:wellposed} (i) $L\circ L=\hbar\,\mathrm{ad}_{\mathcal B}$ at this order. The arity-two component is [HK, cor:c8-homotopy (ii), thm:c12-opalg] with this right side.
\end{proof}

\begin{definition}[The relations to verify]
\label{def:rels}
On cohomology, with $\delta(P_{v_1})=\delta(M_\omega)=0$ in DR (Corollary \ref{cor:exact}):
\begin{itemize}
\item[(R2)] $\{Q,Q\}=P$: $[\delta(\xi),\rho_0(\xi')]+[\rho_0(\xi),\delta(\xi')]=\delta(P_{\Gamma(\xi,\xi')|_V})=0$.
\item[(R3)] $\{Q,S\}\supset D$: $[\delta(\xi),\rho_0(\eta)]+[\rho_0(\xi),\delta(\eta)]=\frac12(\xi,\eta)\,\delta(\Delta)+\delta(R_{\xi\eta})$, $R_{\xi\eta}$ the $\mathfrak{so}(W^\perp)$-part of $[\xi,\eta]$ [HK, prop:c2-special].
\item[(R4)] rank-one $\xi\in S_+\otimes W$: $\{\delta(\xi),\rho_0(\xi)\}=0$, so the twisted class $[Q^\xi_1]$ is defined [HK, thm:c8-twist].
\item[(R5)] $[K,P]\supset D$: $[\delta(K_{v_1}),\rho_0(P_{v_1'})]=(v_1,v_1')\,\delta(\Delta)$.
\item[(R6)] $\{S,S\}=K$: $[\delta(\eta_1),\rho_0(\eta_2)]+[\rho_0(\eta_1),\delta(\eta_2)]=\delta([\eta_1,\eta_2])$.
\item[(R7)] $[D,x]$: $[\delta(\Delta),\rho_0(x)]=0$ on a family at one point graded by the classical weight, for every generator $x$; for $x=\xi$ this is $[D,Q]$: $D_1$ is constant on classical supermultiplets.
\item[(R9)] $[K,Q]=S$: $[\delta(K_{v_1}),\rho_0(\xi)]+[\rho_0(K_{v_1}),\delta(\xi)]=\delta(\rho(v_1)\xi)$; $[P,S]=Q$: $[\rho_0(P_{v_1}),\delta(\eta)]=\delta(\rho(v_1)\eta)$, which is the position decomposition of Proposition \ref{prop:QS}.
\item[(R10)] the Konishi relation: (R3) on the trace bilinear $\sum_I(A_I,A_I)$ and its first descendants, whose right side is $\frac12(\xi,\eta)D_1=-24(\xi,\eta)$ units.
\item[(R11)] joining and lifting: on a non-split extension $0\to N\to E\to C\to0$ with classical joining cocycle $\beta$ and one-loop lifting $\kappa$, $\kappa(\eta)\beta(\xi)c_0=-\gamma(\xi,\eta)c_0$ [FR, Thm 4.9 (iv)]; in the $\mathfrak{su}(2|3)$ sector this is the closure relation $c_9c_{10}=\frac12c_1$ of [Bei03b, (3.36)] between the coefficients of the length-changing supercharges, read from the tree-level BV data, and the one-loop dilatation normalisation $c_1$ [FR, §7].
\end{itemize}
\end{definition}

\begin{theorem}[Exactness table]
\label{thm:table}
In the minimal scheme ($\psi=0$) the relations hold on cohomology as follows; ``MS'' means with $F_1=0$ in the components listed, ``$F$'' with the restoring counterterms of Theorem \ref{thm:restore}, ``test'' means that the vanishing of the defect \eqref{eq:reldefect} with $F_1=0$ is a prediction to be computed.

\mbox{}\par\smallskip\noindent{\small\begin{tabular}{p{0.22\textwidth}|p{0.27\textwidth}|p{0.22\textwidth}|p{0.2\textwidth}}
relation & obstruction components entering & Scheme D & Scheme B\\
\hline
gauge BRST, (R1) & $\mathrm{Ob}^{(0)}$ & MS & MS\\
$P$, $M$ covariance & $\mathrm{Ob}^{(1)}(P),\mathrm{Ob}^{(1)}(M)$ & MS & MS\\
$R$ covariance & $\mathrm{Ob}^{(1)}(R_r)$ & MS & MS on $\mathfrak{so}(4)_J$; $F$ otherwise\\
(R7) along $P,M,R$ & as above & MS & as above\\
(R2) $\{Q,Q\}=P$ & $\mathrm{Ob}^{(1)}(\xi)$, $\mathrm{Ob}^{(2)}(\xi,\xi')$ & test (Fierz only) & $F$ (field counting)\\
(R4) twisted & $\mathrm{Ob}^{(1)}(\xi)$, $\mathrm{Ob}^{(2)}(\xi,\xi)$ & test & $F$\\
(R7) $[D,Q]$ & $\mathrm{Ob}^{(1)}(\Delta)$, $\mathrm{Ob}^{(1)}(\xi)$ & $F$ along $\Delta$ (field counting); test along $\xi$ & $F$\\
(R3) $\{Q,S\}\supset D$, (R10) & $\mathrm{Ob}^{(1)}(\xi,\eta,\Delta)$, $\mathrm{Ob}^{(2)}(\xi,\eta)$ & $F$ (the $(4-D)\mathcal N(\eta)$ and $\Delta$ parts) & $F$\\
(R5) $[K,P]\supset D$ & $\mathrm{Ob}^{(1)}(K),\mathrm{Ob}^{(1)}(\Delta)$ & $F$ along $\Delta,K$ & $F$\\
(R6) $\{S,S\}=K$ & $\mathrm{Ob}^{(1)}(\eta)$, $\mathrm{Ob}^{(2)}(\eta,\eta')$ & $F$ & $F$\\
(R9) & $\mathrm{Ob}^{(1)}(K,\xi,\eta)$, $\mathrm{Ob}^{(2)}(K,\xi)$ & $F$ & $F$\\
(R11) joining/lifting & as (R3) & as (R3) & as (R3)\\
\end{tabular}}\par\smallskip

The entries ``$F$'' along $\Delta$, $K$ and the $(4-D)\mathcal N(\eta)$-part along $\eta$ are not symmetry violations: they carry the anomalous dimensions of the elementary fields and the conformal anomaly of the regulated action, and the corresponding $F_1$ are field-counting functionals and their position moments. The entries ``test'' in Scheme D are the prerequisite and part (a) of Problem \ref{q:Q1}.
\end{theorem}

\begin{proof}
Theorem \ref{thm:rel} with Theorem \ref{thm:restore} and Corollary \ref{cor:exact}. Along $\xi$ in Scheme D, $b_D$ has no explicit $D$-dependence (neither $S^{(1)}_D(\xi)$ nor $\xi.\Psi$ contains $D$), so by Theorem \ref{thm:obs} (i) $\mathrm{Ob}^{(1)}(\xi)=-\mathrm{LIM}\,\partial_D\mathcal R_{\mathfrak F(\xi)}$, a pure Fierz-insertion term; the same holds for $\mathrm{Ob}^{(2)}(\xi,\xi')$ with the Fierz-null $\widehat\Delta^{(2)}(\xi,\xi')$ of Theorem \ref{thm:breakD} (v). In Scheme B, $\mathcal E(\xi)$ contributes already to the two-leg part, as for $\gamma_5$ of 't Hooft--Veltman type in [S05, (4.5)].
\end{proof}

\begin{definition}[Verification protocol]
\label{def:protocol}
For each relation of Definition \ref{def:rels} and each family $\mathfrak F$ (scalar bilinears; the trace bilinear and its descendants of weight $\frac52$ and $3$; the $\mathfrak{su}(2)$ sector of two-letter words; twist two at lowest spin):
\begin{itemize}
\item[1.] compute $\delta(x)$ for the generators involved from \eqref{eq:masterD} (two-body kernels $\Upsilon^D_x$, one-body terms, $\{F^{(1)}_1(x),-\}$);
\item[2.] compute the arity-two kernels $\Upsilon^D_{x,y}$ (cycles through the pair with two arity-one vertices or one arity-two vertex, evanescent vertex once) and $\chi_{x,y}$ of [HK, thm:c12-opalg (ii)];
\item[3.] evaluate the left side of \eqref{eq:reldefect} on a basis of $\mathfrak F$ as an exact matrix identity, subtract the $\Lambda^{(0)}$-exact part $\Lambda^{(0)}\chi_{x,y}$, and compare the remainder with the right side computed from $\mathcal B$;
\item[4.] repeat with $F_1=0$ (MS) and with the restoring $F_1$; record which relations hold in which scheme, and the scheme-independent combinations (classes, spectra).
\end{itemize}
Route A of Section \ref{sec:engines} performs 1--4 with residues; Route B evaluates the same relations on renormalised Green functions with operator insertions, without the quantum action principle.
\end{definition}

\begin{remark}[Two-component form of the relations]
With $Q^i_\alpha$, $\bar Q_{\dot\alpha i}$, $S^\alpha{}_i$, $\bar S^{\dot\alpha i}$ labelled by the components $\xi_\alpha{}^i$, $\bar\xi_{\dot\alpha i}$, $\eta^\alpha{}_i$, $\bar\eta^{\dot\alpha i}$ [HK, def:c8-twocomp] and the brackets of [HK, prop:c2-special]:
\begin{itemize}
\item (R2): the one-loop parts of $\{Q^i_\alpha,\bar Q_{\dot\beta j}\}$ and $\{Q^i_\alpha,Q^j_\beta\}$ vanish on cohomology, their right sides being $\delta(P)=0$; the translation is $\Gamma(\xi,\bar\xi')|_V^m=i\,\bar\xi'_{\dot\alpha i}\bar\sigma^{m\dot\alpha\alpha}\xi_\alpha{}^i$.
\item (R3): the dilatation part of $[\xi,\eta]$ is $\frac12(\xi,\eta)\Delta$ with $(\xi,\eta)=\xi_\alpha{}^i\eta^\alpha{}_i-\bar\xi_{\dot\alpha i}\bar\eta^{\dot\alpha i}$, so the one-loop part of $\{Q^i_\alpha,S^\beta{}_j\}$ contains $\frac12\delta^i_j\delta_\alpha{}^\beta D_1$, plus the $\mathfrak{su}(4)$ generator, whose one-loop part vanishes in Scheme D.
\item (R5): $[K_m,P_n]$ at one loop is $\delta_{mn}D_1$, the rotation part being $M_1=0$.
\item (R6): the one-loop part of $\{S^\alpha{}_i,\bar S^{\dot\beta j}\}$ is $\delta$ of $-K_{v_1}$ with $v_1=\Gamma(\eta,\bar\eta')|_V$.
\end{itemize}
\end{remark}
